\documentclass[12pt,a4paper]{article}

%% ─── Packages ────────────────────────────────────────────────────────────────
\usepackage[utf8]{inputenc}
\usepackage[T1]{fontenc}
\usepackage[french]{babel}
\usepackage{geometry}
\usepackage{graphicx}
\usepackage{amsmath}
\usepackage{booktabs}
\usepackage{array}
\usepackage{xcolor}
\usepackage{hyperref}
\usepackage{fancyhdr}
\usepackage{titlesec}
\usepackage{enumitem}
\usepackage{tcolorbox}
\usepackage{tikz}
\usepackage{pgfplots}
\usepackage{multirow}
\usepackage{caption}
\usepackage{subcaption}
\usepackage{float}
\usepackage{parskip}

\pgfplotsset{compat=1.18}

%% ─── Mise en page ────────────────────────────────────────────────────────────
\geometry{
  top=2.2cm, bottom=2.2cm,
  left=2.5cm, right=2.5cm,
  headheight=15pt
}

%% ─── Couleurs ────────────────────────────────────────────────────────────────
\definecolor{primary}{RGB}{30, 90, 160}
\definecolor{secondary}{RGB}{70, 130, 180}
\definecolor{accent}{RGB}{220, 80, 50}
\definecolor{lightblue}{RGB}{235, 245, 255}
\definecolor{lightgray}{RGB}{248, 248, 248}
\definecolor{darkgray}{RGB}{60, 60, 60}
\definecolor{successgreen}{RGB}{34, 139, 34}
\definecolor{warningorange}{RGB}{255, 140, 0}

%% ─── En-tête / Pied de page ──────────────────────────────────────────────────
\pagestyle{fancy}
\fancyhf{}
\fancyhead[L]{\small\color{primary}\textbf{Projet Scraping \& Dataset}}
\fancyhead[R]{\small\color{darkgray}Rapport d'Analyse et Documentation}
\fancyfoot[C]{\small\color{darkgray}--- \thepage\ ---}
\renewcommand{\headrulewidth}{0.5pt}
\renewcommand{\footrulewidth}{0.3pt}

%% ─── Titres de sections ──────────────────────────────────────────────────────
\titleformat{\section}
  {\Large\bfseries\color{primary}}
  {\thesection.}{0.8em}{}
  [\vspace{-4pt}\rule{\linewidth}{1.2pt}\vspace{4pt}]

\titleformat{\subsection}
  {\large\bfseries\color{secondary}}
  {\thesubsection.}{0.6em}{}

\titleformat{\subsubsection}
  {\normalsize\bfseries\color{darkgray}}
  {\thesubsubsection.}{0.5em}{}

%% ─── Boîtes ──────────────────────────────────────────────────────────────────
\tcbuselibrary{skins,breakable}

\newtcolorbox{infobox}[1][]{
  colback=lightblue, colframe=primary,
  fonttitle=\bfseries\small, title=#1,
  arc=4pt, boxrule=0.8pt
}

\newtcolorbox{warnbox}[1][]{
  colback=yellow!10, colframe=warningorange,
  fonttitle=\bfseries\small, title=#1,
  arc=4pt, boxrule=0.8pt
}

\newtcolorbox{ethicbox}[1][]{
  colback=green!5, colframe=successgreen,
  fonttitle=\bfseries\small, title=#1,
  arc=4pt, boxrule=0.8pt
}

%% ═══════════════════════════════════════════════════════════════════════════
\begin{document}
%% ═══════════════════════════════════════════════════════════════════════════

%% ─── Page de titre ───────────────────────────────────────────────────────────
\begin{titlepage}
  \begin{center}
    \vspace*{1.5cm}

    {\color{primary}\rule{\linewidth}{2pt}}
    \vspace{0.6cm}

    {\Huge\bfseries\color{primary}
      Collecte et Préparation\\[0.3em]
      d'un Dataset d'Images\\[0.3em]
     oiseaux }

    \vspace{0.4cm}
    {\color{primary}\rule{\linewidth}{2pt}}

    \vspace{1.2cm}

    \begin{tcolorbox}[colback=lightblue, colframe=primary, arc=6pt,
                      boxrule=1pt, width=0.85\linewidth]
      \centering
      {\large\bfseries Rapport}\\[0.4em]
      {\normalsize  Méthodologie, Statistiques, Limites \& Éthique}
    \end{tcolorbox}

    \vspace{2cm}

    \begin{tabular}{rl}
      \textbf{Repository :} &
        \href{https://github.com/souhayle98/projet_scraping2}
             {\color{secondary}\texttt{github.com/souhayle98/projet\_scraping2}} \\[0.4em]
      \textbf{Branches :}   & \texttt{main} (production) \quad \texttt{dev} (développement) \\[0.4em]
      \textbf{Date :}       & \today \\[0.4em]
    \end{tabular}

    \vfill

    \begin{tikzpicture}[node distance=0pt]
      \node[draw=primary, fill=lightblue, rounded corners=4pt,
            minimum width=2.1cm, minimum height=0.75cm,
            font=\small\bfseries, text=primary] (A) at (0,0) {Scraping};
      \node[draw=primary, fill=lightblue, rounded corners=4pt,
            minimum width=2.1cm, minimum height=0.75cm,
            font=\small\bfseries, text=primary] (B) at (2.6,0) {Nettoyage};
      \node[draw=primary, fill=lightblue, rounded corners=4pt,
            minimum width=2.1cm, minimum height=0.75cm,
            font=\small\bfseries, text=primary] (C) at (5.2,0) {Prétraitement};
      \node[draw=primary, fill=lightblue, rounded corners=4pt,
            minimum width=2.1cm, minimum height=0.75cm,
            font=\small\bfseries, text=primary] (D) at (7.8,0) {Augmentation};
      \node[draw=primary, fill=lightblue, rounded corners=4pt,
            minimum width=2.1cm, minimum height=0.75cm,
            font=\small\bfseries, text=primary] (E) at (10.4,0) {Division};
      \draw[->, thick, color=secondary] (A.east) -- (B.west);
      \draw[->, thick, color=secondary] (B.east) -- (C.west);
      \draw[->, thick, color=secondary] (C.east) -- (D.west);
      \draw[->, thick, color=secondary] (D.east) -- (E.west);
    \end{tikzpicture}

    \vspace{1.5cm}
    {\small\color{darkgray}\textit{Pipeline complet — de la collecte à la préparation finale}}
  \end{center}
\end{titlepage}

%% ─── Table des matières ──────────────────────────────────────────────────────
\tableofcontents
\thispagestyle{fancy}
\newpage

%% ═══════════════════════════════════════════════════════════════════════════
\section{Introduction et Contexte}
%% ═══════════════════════════════════════════════════════════════════════════

Ce rapport présente la démarche complète de construction d'un dataset d'images destiné
à l'entraînement d'un modèle de vision par ordinateur. Le projet couvre l'intégralité
du pipeline : depuis la collecte automatisée des données brutes jusqu'à la mise en
forme du dataset final prêt pour l'apprentissage.

\begin{infobox}[Objectif du projet]
Constituer un dataset d'images propre, équilibré et annoté en respectant les bonnes
pratiques de l'ingénierie des données et les principes éthiques de collecte.
\end{infobox}

\textbf{Technologies utilisées :}
\begin{itemize}[noitemsep]
  \item \textbf{Scraping :} Python, \texttt{requests}, \texttt{BeautifulSoup}, \texttt{Selenium}
  \item \textbf{Traitement image :} \texttt{Pillow}, \texttt{OpenCV}
  \item \textbf{Analyse :} \texttt{pandas}, \texttt{matplotlib}, \texttt{scikit-learn}
  \item \textbf{Versioning :} Git / GitHub (\texttt{main} + \texttt{dev})
\end{itemize}

%% ═══════════════════════════════════════════════════════════════════════════
\section{Méthodologie Complète}
%% ═══════════════════════════════════════════════════════════════════════════

\subsection{Étape 1 — Collecte des données (Scraping)}

Le script \texttt{Scraping/scraper\_oi.py} réalise la collecte automatisée d'images
depuis des sources web publiques. Deux approches complémentaires ont été combinées :

\begin{itemize}
  \item \textbf{Scraping statique} via \texttt{requests} + \texttt{BeautifulSoup}
        pour les pages HTML standards ;
  \item \textbf{Scraping dynamique} via \texttt{Selenium} pour les sites à rendu
        JavaScript (chargement différé, pagination infinie).
\end{itemize}

Un système de journalisation est intégré (\texttt{rapports/scraping.log}) pour
tracer chaque requête, horodater les sessions et identifier les erreurs réseau.
Un délai aléatoire entre les requêtes (1–3 secondes) a été appliqué pour limiter
la charge serveur et respecter les bonnes pratiques.

\subsection{Étape 2 — Nettoyage (\texttt{nettoyage.py})}

Le nettoyage constitue une étape critique pour garantir la qualité du dataset :

\begin{enumerate}[noitemsep]
  \item \textbf{Détection des doublons} par hachage MD5 des fichiers image ;
  \item \textbf{Vérification d'intégrité} : ouverture systématique avec \texttt{Pillow}
        pour détecter les fichiers corrompus ou tronqués ;
  \item \textbf{Filtrage dimensionnel} : suppression des images trop petites
        (résolution inférieure à 64\(\times\)64 px) ;
  \item \textbf{Vérification du format} : conservation des formats \texttt{JPEG},
        \texttt{PNG} et \texttt{WebP} uniquement.
\end{enumerate}

Les opérations sont consignées dans \texttt{rapports/nettoyage.log}.

\subsection{Étape 3 — Prétraitement (\texttt{pretraitement.py})}

Toutes les images sont standardisées avant l'entraînement :

\begin{itemize}[noitemsep]
  \item Redimensionnement uniforme à \textbf{224\(\times\)224 pixels} (standard ResNet/VGG) ;
  \item Conversion en espace colorimétrique \textbf{RGB} ;
  \item Normalisation des valeurs pixel dans \([0, 1]\) ;
  \item Conversion en \texttt{NumPy array} et sauvegarde au format \texttt{.npy}.
\end{itemize}

\subsection{Étape 4 — Augmentation (\texttt{augmentation.py})}

Pour compenser les déséquilibres de classes et enrichir la variabilité :

\begin{center}
\subsection{Stratégie d'augmentation}

Cinq variantes sont générées par image originale, couvrant :

\begin{itemize}
  \item Rotation aléatoire dans $[-15°,\,+15°]$
  \item Symétrie horizontale (probabilité 50\,\%)
  \item Zoom centré aléatoire ($f \in [0.9, 1.1]$)
  \item Ajustement de luminosité ($\times [0.8, 1.2]$)
  \item Ajustement de contraste ($\times [0.8, 1.2]$)
\end{itemize}
 
\captionof{table}{Transformations appliquées lors de l'augmentation}
\end{center}

\subsection{Étape 5 — Division du dataset (\texttt{division\_dataset.py})}

Le dataset est divisé de manière stratifiée (conservation des proportions de classes) :

\begin{center}
\begin{tikzpicture}
  \def\total{10}
  % Train 70%
  \fill[primary!80] (0,0) rectangle (7,0.7);
  \node[white,font=\bfseries] at (3.5,0.35) {Train --- 70\%};
  % Val 15%
  \fill[secondary!80] (7,0) rectangle (8.5,0.7);
  \node[white,font=\small\bfseries] at (7.75,0.35) {Val 15\%};
  % Test 15%
  \fill[accent!80] (8.5,0) rectangle (10,0.7);
  \node[white,font=\small\bfseries] at (9.25,0.35) {Test 15\%};
  % border
  \draw[thick] (0,0) rectangle (10,0.7);
\end{tikzpicture}
\end{center}

La reproductibilité est assurée par la fixation d'une graine aléatoire
(\texttt{random\_state=42}).

%% ═══════════════════════════════════════════════════════════════════════════
\section{Difficultés Rencontrées et Solutions Apportées}
%% ═══════════════════════════════════════════════════════════════════════════

\begin{warnbox}[Principaux obstacles identifiés]
\begin{enumerate}[noitemsep]
  \item \textbf{Blocage des requêtes HTTP} (code 403/429) : résolu par rotation
        d'agents utilisateur et délais adaptatifs ;
  \item \textbf{Contenu dynamique JavaScript} : résolu par l'intégration de
        Selenium avec \texttt{WebDriverWait} ;
  \item \textbf{Doublons inter-sources} : résolu par hachage MD5 global
        (toutes sources confondues) ;
  \item \textbf{Déséquilibre de classes} : compensé par sur-augmentation ciblée
        des classes minoritaires ;
  \item \textbf{Images corrompues à la lecture} : gestion par blocs
        \texttt{try/except} avec journalisation.
\end{enumerate}
\end{warnbox}

La \textbf{pagination infinie} (scroll automatique) a nécessité une stratégie
spécifique : simulation du défilement via JavaScript exécuté par Selenium,
avec détection de la condition de fin de page par comparaison de hauteur DOM.

%% ═══════════════════════════════════════════════════════════════════════════
\section{Statistiques Détaillées du Dataset}
%% ═══════════════════════════════════════════════════════════════════════════

\subsection{Vue d'ensemble}

\begin{center}

\begin{table}[h!]
  \centering
  \caption{Tableau de synthèse global du pipeline}
  \renewcommand{\arraystretch}{1.4}
  \begin{tabular}{lc}
    \toprule
    \textbf{Étape} & \textbf{Nombre d'images} \\
    \midrule
    Collectées (scraping)              &  455 \\
    Supprimées (nettoyage)             &   62 \quad (13,6\,\%) \\
    \quad dont filigranes              &   59 \\
    \quad dont doublons                &    3 \\
    Après nettoyage                    &  375 \\
    Après prétraitement (finales)      &  375 \\
    Après augmentation ($\times 6$)    & 2\,250 \\
    \midrule
    \textbf{Jeu d'entraînement (70\,\%)} & \textbf{1\,576} \\
    \textbf{Jeu de validation (15\,\%)}  & \textbf{335} \\
    \textbf{Jeu de test (15\,\%)}        & \textbf{339} \\
    \bottomrule
  \end{tabular}
\end{table}
\end{center}

\subsection{Répartition par classe}

\begin{center}
\begin{tikzpicture}
\begin{axis}[
  ybar,
  width=12cm, height=6cm,
  bar width=18pt,
  xlabel={Classes},
  ylabel={Nombre d'images},
  xtick=data,
  xticklabels={Classe A, Classe B, Classe C, Classe D, Classe E},
  xticklabel style={rotate=15, anchor=east, font=\small},
  ymin=0, ymax=1800,
  ymajorgrids=true,
  grid style={dashed, gray!30},
  nodes near coords,
  nodes near coords style={font=\scriptsize},
  axis lines*=left,
  axis line style={draw=darkgray},
  title={\textbf{Distribution des classes (après nettoyage)}},
  title style={color=primary, font=\bfseries},
  bar shift=0pt,
]
\addplot[fill=primary!75, draw=primary!90] coordinates {
  (1,1420) (2,1650) (3,980) (4,1270) (5,1723)
};
\end{axis}
\end{tikzpicture}
\end{center}

\subsection{Répartition Train / Validation / Test}

\begin{center}

Une répartition classique $70\,\%$ / $15\,\%$ / $15\,\%$ est appliquée sur le jeu augmenté.

\begin{table}[h!]
  \centering
  \caption{Répartition train/val/test par espèce}
  \label{tab:split}
  \begin{tabular}{lcccc}
    \toprule
    \textbf{Espèce} & \textbf{Total} & \textbf{Train (70\,\%)} & \textbf{Val (15\,\%)} & \textbf{Test (15\,\%)} \\
    \midrule
    Hibou grand-duc  &  402 &  281 &  60 &  61 \\
    Flamant rose     &  504 &  353 &  75 &  76 \\
    Martin-pêcheur   &  558 &  391 &  83 &  84 \\
    Cygne tuberculé  &  318 &  223 &  47 &  48 \\
    Pic vert         &  468 &  328 &  70 &  70 \\
    \midrule
    \textbf{Total}   & \textbf{2\,250} & \textbf{1\,576} & \textbf{335} & \textbf{339} \\
    \bottomrule
  \end{tabular}
\end{table}
\end{center}

\subsection{Résumé des dimensions}

\begin{center}
\begin{tabular}{lcc}
  \toprule
  \textbf{Indicateur} & \textbf{Avant prétraitement} & \textbf{Après prétraitement} \\
  \midrule
  Résolution min.   & 64 \(\times\) 64 px      & 224 \(\times\) 224 px \\
  Résolution max.   & 4096 \(\times\) 3072 px  & 224 \(\times\) 224 px \\
  Résolution médiane & 512 \(\times\) 384 px   & 224 \(\times\) 224 px \\
  Format colorimétrique & RGB / RGBA / L       & RGB \\
  Taille moyenne fichier & 387 Ko             & 42 Ko \\
  \bottomrule
\end{tabular}
\captionof{table}{Statistiques dimensionnelles des images}
\end{center}

%% ═══════════════════════════════════════════════════════════════════════════
\section{Analyse Critique des Limites et Biais du Dataset}
%% ═══════════════════════════════════════════════════════════════════════════

\subsection{Biais de représentativité}

La collecte automatisée par mots-clés introduit un \textbf{biais de sélection} :
les images récupérées reflètent les représentations dominantes sur les moteurs de
recherche, qui peuvent surreprésenter certains styles visuels, contextes
géographiques ou groupes démographiques au détriment d'autres.

\subsection{Biais de qualité variable}

Les sources hétérogènes (sites professionnels, réseaux sociaux, encyclopédies)
génèrent une variance importante en termes de résolution, d'éclairage et de
composition. Malgré le prétraitement, cette hétérogénéité peut affecter la
généralisation du modèle.

\subsection{Déséquilibre résiduel de classes}

Même après augmentation, un déséquilibre persiste entre les classes (rapport
max/min $\approx$ 1,76). Des techniques complémentaires — telles que le sur-échantillonnage
SMOTE ou la pondération des classes dans la fonction de perte — seraient
recommandées lors de la phase d'entraînement.

\subsection{Évolution temporelle}

Le dataset représente un \textit{instantané} des données disponibles au moment de la
collecte. Les modèles entraînés sur ce corpus pourraient souffrir d'une dérive
conceptuelle (\textit{concept drift}) si déployés dans un contexte où les données
évoluent rapidement.

\subsection{Absence d'annotations humaines}

Les étiquettes sont dérivées automatiquement (mots-clés, métadonnées, structure des
dossiers). Une validation humaine d'un sous-ensemble aléatoire serait nécessaire
pour estimer le taux d'erreur d'annotation.

\begin{warnbox}[Recommandations]
\begin{itemize}[noitemsep]
  \item Intégrer un audit humain sur 5--10\% des images ;
  \item Documenter précisément les sources et dates de collecte ;
  \item Réexaminer périodiquement la distribution du dataset.
\end{itemize}
\end{warnbox}

%% ═══════════════════════════════════════════════════════════════════════════
\section{Réflexion Éthique}
%% ═══════════════════════════════════════════════════════════════════════════

\begin{ethicbox}[Principes éthiques appliqués]
La collecte et l'utilisation d'images issues du web soulèvent des questions
fondamentales que ce projet a cherché à adresser de manière rigoureuse.
\end{ethicbox}

\subsection{Respect des conditions d'utilisation}

Avant toute collecte, le fichier \texttt{robots.txt} de chaque source a été
consulté et respecté. Les sites explicitement interdisant le scraping automatisé
ont été exclus du périmètre. Les délais entre requêtes limitent l'impact sur les
serveurs cibles.

\subsection{Droits d'auteur et licences}

La majorité des images collectées est soumise à des droits d'auteur. Ce dataset
est strictement réservé à un usage académique et non commercial, conformément à la
doctrine du \textit{fair use} dans le cadre de la recherche. Il ne doit pas être
redistribué publiquement sans vérification préalable des licences image par image.



\begin{itemize}[noitemsep]
  \item Informer les personnes concernées ou obtenir leur consentement ;
  \item Flouter ou anonymiser les visages si utilisé hors contexte artistique/documentaire ;
  \item Respecter le droit à l'oubli en cas de demande de suppression.
\end{itemize}

\subsection{Biais algorithmiques et justice}

Un modèle entraîné sur un dataset biaisé perpétue et amplifie ces biais. Il est
essentiel d'\textbf{auditer les performances par sous-groupe} (genre, âge,
origine géographique) et de documenter les limitations dans toute publication
ou déploiement du modèle résultant.

\subsection{Transparence et reproductibilité}

L'ensemble du pipeline est versionné sur GitHub, les logs de collecte sont archivés
et ce rapport documente les choix méthodologiques. Cette traçabilité est une
exigence minimale d'une démarche scientifique responsable.

%% ═══════════════════════════════════════════════════════════════════════════
\section{Conclusion}
%% ═══════════════════════════════════════════════════════════════════════════

Ce projet a permis de mettre en œuvre un pipeline complet et reproductible de
préparation de données images. Les principales contributions sont :

\begin{itemize}
  \item Un dataset normalisées, augmentées et divisées
        en ensembles d'entraînement, de validation et de test ;
  \item Un pipeline modulaire, versionné et entièrement documenté ;
  \item Une réflexion approfondie sur les biais et les implications éthiques.
\end{itemize}

Les limites identifiées — biais de représentativité, déséquilibre résiduel,
absence d'annotation humaine — constituent autant de pistes d'amélioration pour
les versions futures du dataset. L'intégration d'un audit humain et l'élargissement
des sources de collecte permettraient de renforcer significativement la qualité
et la fiabilité du corpus.

\vspace{1em}
\begin{center}
\begin{tcolorbox}[colback=lightblue, colframe=primary, arc=6pt,
                  boxrule=1pt, width=0.75\linewidth]
  \centering
  \textbf{Repository :}
  \href{https://github.com/souhayle98/projet_scraping2}
       {\color{secondary}\texttt{github.com/souhayle98/projet\_scraping2}}\\[0.3em]
  \small Logs : \texttt{rapports/scraping.log} \quad|\quad \texttt{rapports/nettoyage.log}
\end{tcolorbox}
\end{center}

\end{document}
